% !TeX root = ../教材题库_全部条目_转写工作稿.tex
% Source: XJ-B1-C02-PR-P042-Y0646
% 第2章 / 2.1 相等关系与不等关系 / 教材 p.42 / PDF 第49页
% OCR未校对：公式与文字均可能有误；未完成数学TeX转写；禁止直接发给学生
% 原题及相邻原解截图（核对用）：../题目原图/XJ-B1-C02-PR-P042-Y0646-01.webp
% 本文件只在转写完成后作为题干使用；OCR 条目尚不是可靠题干。
\begin{bankitem}{p.42 练习 第1题}{第2章 / 2.1 相等关系与不等关系 / 教材 p.42 / PDF 第49页}{XJ-B1-C02-PR-P042-Y0646}{OCR未校对：公式与文字均可能有误；未完成数学TeX转写；禁止直接发给学生}
\begin{ocrdraft}
[教材p.42 / PDF第49页]
\par
1.甲、乙两同学分别解“设e0，十亠），求二次函数3，一2T2十1的最小
\par
值”的过程如下：
\par
甲：3，一2十1冫22
\par
1一2x，
\par
又冫1，所以22．
\par
从而3，冫2、／冫2，即，的最小值是2．
\par
乙：因为函数3，一2十1在区间0，十。。）上的图象随着增大而逐渐上升，
\par
即y随增大而增大，所以3，的最小值是2\ensuremath{\times}12十1一3．
\par
试判断谁错，错在何处？
\end{ocrdraft}
\end{bankitem}
